\documentclass[../../../../main.tex]{subfiles}
\begin{document}

\begin{flushright}
\footnotesize\textit{【自動文字起こし・要確認】}
\end{flushright}

[解] (1) $T: ax+y+z+a-2=0$ とする。$T$の法線ベクトルの1つに$\vec{n} = \begin{pmatrix} a \\ 1 \\ 1 \end{pmatrix}$がある。題意の点$P$とすると、$OP \perp T$ だから、$k \in \mathbb{R}$として
\begin{equation*}
\overrightarrow{OP} = k\vec{n} \quad \cdots \text{\textcircled{1}}
\end{equation*}
とかける。これが$T$上の点だから、
\begin{align*}
a(ka) + (k) + (k) + a-2 &= 0 \qquad \therefore k = \frac{2-a}{a^2+2}
\end{align*}
となって \textcircled{1}から
\begin{align*}
P\left( \frac{(2-a)a}{a^2+2}, \frac{2-a}{a^2+2}, \frac{2-a}{a^2+2} \right) \quad \text{\#\ \ }
\end{align*}

(2) 平面と点のキョリ公式から
\begin{align*}
|OP| &= \frac{|a-2|}{\sqrt{a^2+1+1}} = \frac{|a-2|}{\sqrt{a^2+2}} \quad (\equiv f(a) \text{とする})
\end{align*}
である。まず $\left\{ f(a) \right\}^2 = \frac{(a-2)^2}{a^2+2}$ の値域をもとめる。
\begin{align*}
\left\{ f(a)^2 \right\}' &= \frac{2(a-2)(a^2+2) - 2a(a-2)^2}{(a^2+2)^2} = \frac{2(a-2)(2a+2)}{(a^2+2)^2}
\end{align*}
から増減表をうる。

\begin{center}
\begin{tabular}{|c|c|c|c|c|c|}
\hline
$a$ & $\cdots$ & $-1$ & $\cdots$ & $2$ & $\cdots$ \\ \hline
$f'$ & $+$ & $0$ & $-$ & $0$ & $+$ \\ \hline
$f^2$ & $\nearrow$ & $3$ & $\searrow$ & $0$ & $\nearrow$ \\ \hline
\end{tabular}
\end{center}

したがって、$0 \le f(a)$とあわせて、$\max f(a) = \sqrt{3}$である。$V(a) \text{が}\min \iff f(a) \text{が}\max$ だから\\
$a=-1$の時、
\begin{align*}
\min V(a) &= \pi \int_{\sqrt{3}}^3 (9-x^2)dx = \pi \left[ 9x - \frac{1}{3}x^3 \right]_{\sqrt{3}}^3 = \pi (18 - 8\sqrt{3}) \quad \text{\#\ \ }
\end{align*}
をとる。

\end{document}
